%% LyX 2.4.4 created this file.  For more info, see https://www.lyx.org/.
%% Do not edit unless you really know what you are doing.
\documentclass[english]{article}
\usepackage[T1]{fontenc}
\usepackage[latin9]{inputenc}
\usepackage{enumitem}
\usepackage{amsmath}
\usepackage{amsthm}
\usepackage{geometry}
\geometry{verbose,tmargin=1in,bmargin=1in,lmargin=1in,rmargin=1in}

\makeatletter
%%%%%%%%%%%%%%%%%%%%%%%%%%%%%% Textclass specific LaTeX commands.
\numberwithin{equation}{section}
\numberwithin{figure}{section}
\newlength{\lyxlabelwidth}      % auxiliary length 

%%%%%%%%%%%%%%%%%%%%%%%%%%%%%% User specified LaTeX commands.
\usepackage{fancyhdr}
\usepackage{lastpage}
\pagestyle{fancy}
\usepackage{enumitem}
% Added by lyx2lyx
\setlength{\parskip}{\smallskipamount}
\setlength{\parindent}{0pt}

\makeatother

\usepackage{babel}
\begin{document}
\lhead{MTH 331(A) James Ashe\\ 
\textbf{Quiz} September, 26 2012}
\chead{\textbf{Name:}}
\cfoot{\thepage{}~of~\pageref{LastPage}}
%%
%%\setenumerate[0]{leftmargin=12pt,itemindent=0pt,labelwidth=0pt}
\renewcommand{\labelenumi}{\textbf{\arabic{enumi}.}}
\renewcommand{\labelenumii}{\textbf{\alph{enumii}.}}

\vspace*{-.75cm}

\textbf{Justify your answers and show your work.}

\emph{Find the limit of the following sequences or determine that
the limit does not exist.}
\begin{enumerate}[resume]
\item $\left\{ \frac{n^{10}}{3n^{10}+4}\right\} $ \vspace{.75in}
\item $\left\{ \frac{\cos\pi n}{\sqrt{n}}\right\} $\vspace{.75in}
\item ${\displaystyle \left\{ \frac{n!}{n^{n}}\right\} }$\vspace{.75in}
\end{enumerate}
\emph{Evaluate the geometric series or explain why it diverges. }
\begin{enumerate}[resume]
\item ${\displaystyle \sum_{k=0}^{\infty}\left(\frac{3}{5}\right)^{k}}$\vspace{.75in}
\end{enumerate}
\emph{Consider the following infinite series.}
\begin{enumerate}[resume]
\item ${\displaystyle \sum_{k=1}^{\infty}\frac{1}{4k^{2}-1}}$.

\begin{enumerate}
\item Write out the first four terms of the sequence of partial sums.\vspace{1.25in}
\item Use the results of part (a) to propose an explicit formula for $S_{n}$.\vspace{.75in}
\end{enumerate}
\item [\textbf{bonus.}]Prove that ${\displaystyle \sum_{k=1}^{\infty}2k-1=n^{2}}$. 
\end{enumerate}

\end{document}
